\honest{Ethical Injunctions}{Eliezer Yudkowsky}{2008}

I am professionally interested in things you should not do even when they seem right. Take an AI whose goals are still unfinished. If it ever decides that fooling its programmers is the right thing to do, that probably means its goals have gone wrong, so it should shut itself down at once, without recalculating. I think this ``should'' be possible, but I do not know the mathematics, and for now it is ``only a dream.'' \nb{In 2015 Yudkowsky and colleagues reported that no proposal for an AI that shuts down safely when a button is pressed yet met all their requirements, and called the problem ``wide-open.''}

The human version is an ethical injunction: never murder an innocent person who has helped you, even if it seems right, because it is far more likely that you have made a mistake. Is that reasonable? In 1944 Norwegian saboteurs sank the ferry Hydro to stop a shipment of heavy water to Germany. Knut Haukelid did not warn the watchman who had let them stay aboard, and eighteen people died. ``And that was, effectively, the end of the Nazi atomic weapons program.'' \nb{The details that can be checked are right, but the program went on: its last reactor experiment, in early 1945, had 1.5 tons of heavy water, and the ferry carried about 500 kg.} Germany very likely would not have got the Bomb anyway. I cannot say a word against it.

For self-deception I know of no exception. People have often been better off with a false belief, just as some lottery ticket always wins, but no one can know in advance which. Self-deceptions are ``the worst kind of black swan bets'': one blowup, such as praying instead of seeing a doctor about a lump, undoes all the good. All the happiness the thought of an afterlife ever gave, I say, ``has now been more than cancelled'' by humanity's failure to adopt cryonic preservation. \nb{Neither quantity is measured, and the claim assumes that cryonics would have worked, which has not been shown.}

So injunctions have two grounds: corrupted hardware, which makes corruption likelier than a correct calculation in some kinds of case, and black swans, actions that blow up for reasons outside the decider's model. Then I state the obvious objection: why not simply redo the calculation with these reasons included? \nb{This is the standard objection to rule-based consequentialism, and the post states it fairly.} I give three replies. It is the cleverness I warn against. Meta-reasoning can be corrupted too, so an AI that has decided on paperclips should shut down, not reassess. And my ethics have protected me from myself before. \nb{That history is told in ``Protected From Myself,'' which is not in the book.} Still, I grant that ``Am I still dumber than my ethics?'' does not automatically get the answer yes.

My own revisions have made my ethics stricter; other people's go the other way. Those who urged me to lie about the Singularity ``seem to have no idea of the risks,'' and I am ``reasonably sure'' this is because they have ``no clue.'' To those who think themselves smarter than their ethics: ``Ha.'' To those who learned all this from my sequence and now think they will take it into account: ``Double ha.'' People who excuse themselves, in my experience, move ``Always'' toward lenience. \nb{Mill made a similar observation in 1861: every creed allows some latitude for circumstances, and there ``self-deception and dishonest casuistry get in.''} As Hobbes said, their price is so low. \nb{This Hobbes is the tiger in Bill Watterson's comic strip \textsc{Calvin and Hobbes}.}

I cannot endorse absolute injunctions, though the one against self-deception has ``tremendous force.'' I end by declaring myself ``completely unimpressed'' with those who have told me that an unethical act is rational.
